\chapter{Introdução}
O reconhecimento facial tem sido um instrumento amplamente utilizado para contextos de identificação e garantia de segurança. O mapeamento de padrões e características é essencial para o bom funcionamento de um sistema de proteção. Nesse cenário, o desenvolvimento de interfaces homem-máquina (IHM) tem avançado com o aumento da capacidade computacional em sistemas embarcados.

Para traduzir essas características em dados quantizados, as Redes Neurais Convolucionais (CNNs) têm se mostrado altamente eficazes. Em uma aplicação real voltada para a autorização da entrada de um indivíduo cadastrado em um estabelecimento, o fluxo de dados exige a aquisição de imagem em tempo real, o pré-processamento e a inferência no modelo de \textit{Deep Learning}. Contudo, a etapa de inferência impõe um alto custo computacional. Processadores de uso geral e microcontroladores processam as matrizes da rede neural de forma puramente sequencial, o que pode representar um gargalo crítico para a percepção de tempo real do usuário.

Para a implementação da rede neural no FPGA, optou-se pelo desenvolvimento manual dos módulos de \textit{hardware} utilizando a Linguagem de Descrição de \textit{Hardware} (HDL) Verilog. Essa abordagem permite maior controle sobre a arquitetura implementada, possibilitando a exploração explícita do paralelismo inerente às operações convolucionais e de multiplicação-acumulação (MAC) presentes na CNN. Além disso, a implementação manual permite avaliar de maneira mais precisa os benefícios e os custos associados ao desenvolvimento de aceleradores dedicados para aplicações de Inteligência Artificial embarcada.

\section{Objetivos}
Este trabalho tem como objetivo o desenvolvimento de uma \textit{Tiny-CNN}, seu respectivo treinamento utilizando a linguagem de programação Python através de um \textit{dataset} com fotos dos integrantes convertidas em escala cinza, a implementação de seu modelo matemático em linguagem de \textit{hardware} Verilog --- além da integração de periféricos, como o protocolo UART e a transmissão de imagem via VGA --- e o teste prático do modelo em uma placa FPGA DE2-115.

\section{Metodologia}
Os dados de treinamento foram obtidos através de uma câmera de \textit{notebook}, em sala de aula, em gravações de vídeos de cada integrante, extraindo os \textit{frames} dos vídeos utilizando a biblioteca openCV em Python e tratando-os em escala cinza e o respectivo treinamento foi realizado mediante Keras. Já o teste de funcionamento do modelo matemático em Verilog baseado na construção do modelo Python foi feito pelo \textit{software} ModelSim, e seu respectivo resultado de simulação comparado com o resultado do modelo Python.
O procedimento adotado foi a construção aditiva e gradual da rede em paralelo com implementação e testes de módulos lógicos em Verilog, orquestrados por uma Máquina de Estados Finitos (FSM) confeccionada também paralelamente.

\section{Justificativa}
O trabalho traz contribuições importantes para a literatura ao fazer um estudo comparativo do processo de inferência de redes neurais convolucionais em diferentes arquiteturas. Além disso, é capaz de atuar como referência em possíveis futuras estruturações de fluxos completos de uma CNN em um FPGA.
\nocite{Hoxha:2025:FAN}
\nocite{MADAOUI2023104927}